% ---------------------------------------------------------------------
%  Abstract
% ---------------------------------------------------------------------
\chapter*{Abstract}
\addcontentsline{toc}{chapter}{Abstract}

Rehosting---executing unmodified embedded firmware inside a full-system
emulator---has become a standard technique for testing, debugging, and
security analysis of embedded software. Frameworks such as Renode
faithfully reproduce the \emph{functional} behaviour of microcontroller
platforms, but their \emph{temporal} behaviour can deviate substantially
from real hardware: emulated clocks advance in coarse scheduling quanta,
interrupt latency is idealised, and physical effects such as oscillator
drift and bus propagation delay are absent entirely. For timing-sensitive
firmware, these deviations can hide real bugs and create false confidence
in test results.

This thesis quantifies and mitigates the temporal deviation between a
physical microcontroller and its rehosted counterpart. Using an
STM32F103RB (Arm Cortex-M3, \SI{72}{\mega\hertz}) running Zephyr RTOS and
an instrumented GPIO-marker firmware, we measure ten timing scenarios on
real hardware with a \SI{100}{\mega\sample\per\second} logic analyser and compare them
against the identical firmware running in Renode~1.15.3. The comparison
reveals systematic deviations of up to three orders of magnitude in
timing \emph{variance}: the emulator is simultaneously too deterministic
for tick-quantised sleeps (jitter ratio ${\approx}\,0.2$) and too noisy
for hardware-timer callbacks (jitter ratio up to $23{\times}$).

Guided by a taxonomy of five deviation sources, we design and implement
the \emph{Temporal Fidelity Layer} (TFL): three composable C\# peripheral
models for Renode---a clock model injecting configurable drift and
Gaussian jitter, a deadline monitor recording task wakeup deviation, and
a bus delay model reproducing inter-node propagation latency. TFL
requires no changes to Renode itself and no changes to the firmware
logic; it is enabled purely through platform description files.

Across seven statistically evaluated scenarios, TFL reduces the
geometric-mean jitter-ratio error from $7.6{\times}$ to $2.9{\times}$
and brings four of seven scenarios within a factor of $1.7$ of hardware
variance. We further characterise the fidelity--performance trade-off of
Renode's scheduling quantum, showing that quantum reduction alone cannot
reach hardware-level fidelity at acceptable simulation cost. All
measurements, models, and analysis scripts are released as open
research artefacts.
